\scanpage{107}
% archive foliation: 106
\deleted{
\[
\underline{\tau_{14}},\tau_{43}\text{ ou }\tau_{46}\ ?\qquad
\tau_{13}\text{ ou }\tau_{16}
\]
\[1\leq i\leq5\qquad\rho_s^i\sigma\rho_s^j\sigma\ldots\rho_s^h\sigma\]
\[\rho_s^{\unsure3}=(\tau_{4,5})\qquad\rho_s^4=(1,2,3)(4)(5)(6)\]
}
\deleted{$(fgh)'=(fg)'h+(fg)h'=\unsure{f'gh|fg'h}+fg'h+fgh'$}

\[a=\underline0\quad\underline b=1\qquad\unsure{s}=\underline\infty\]
\uline{reste}\quad$c,x,y,u,v$
\[f(z)=\lambda\uuline{z^3(z-1)^2(z-c)}\]
\[
\left\{\begin{array}{ll}\underline0&\text{zéro triple}\\\underline1&\text{zéro double}\\c&\text{zéro simple}\end{array}\right.
\qquad\left\{\infty\text{ pôle d'ordre }6\right.
\]
\[
\begin{aligned}
f'(z)&=\lambda\deleted{\unsure{f}}\bigl[z^2(z-1)(\underbrace{3(z-1)(z-c)+2z(z-c)}_{(5z-3)(z-c)}+z(z-1))\bigr]\\
&=\lambda z^2(z-1)[6z^2-(5c+4)z+3c]
\end{aligned}
\]
Ses zéros distincts de $0,1,\infty$ sont les zéros de
\[6z^2-(5c+4)z+3c=0\]
\deleted{\unsure{soient}} $u,v$. Il faut donc \deleted{\unsure{écrire} $f$}
\deleted{\illegible{1}\unsure{que}}
\[
\left\{\begin{array}{l}u,v\text{ distincts entre eux et }\ne0,1,c\\f(u)=f(v)=1\end{array}\right.
\qquad\text{(racines doubles de cette équation}
\]
\[f(z)-1=\deleted{\unsure{\lambda z^2}}\ \lambda(z-u)^2(z-v)^2(z-x)(z-y)\]
